\chapter*{Scoring Functions in Molecular Docking}


Scoring functions are designed to reproduce the binding thermodynamics, \textit{i.e.}, binding enthalpy and entropy (see Eq \textbf{\ref{eq:free-energy}}).

\begin{equation} 
    \label{eq:free-energy}
    \   {\Delta}G_{binding} = {\Delta}H - T{\Delta}S   
\end{equation}

Where  \({\Delta}H\) is the enthalpy component of binding and \(T{\Delta}S\) the entropy component of binding at temperature "$T$".
The enthalpy of binding can be estimated by the number and type of interactions between the receptor and ligand at the atomistic level. Hence, scoring functions often estimate the enthalpy component by summing all interactions of different types. However, this is a non-trivial task, which has also been criticised for being treated as purely additive\cite{bib61, bib62}. Nevertheless, different types of scoring functions employ different methods to fit experimental binding data to develop scoring functions. This leads to different types of scoring functions that can be classified into four types. We briefly discuss the scoring functions and their types, for more detailed discussions the readers are directed to consults these papers\cite{bib51, bib52}.

For a docking tool to be effective and yield reliable result, the search algorithm and the scoring function are crucial and carefully developed.  The single most important question that is often asked about a docking algorithm is “How efficiently can the algorithm differentiate between binders and non-binders?” For the docking algorithm to efficiently differentiate between them, both the conformation search algorithm and the scoring functions must work in tandem at their best capacity. A huge data set of chemical decoys is employed to validate the accuracy, which is spiked with known binders. Chemical decoys are chemical structures that physically resemble known binders but have a different topological arrangements of atoms. The accuracy of the docking algorithm is assessed by its ability to recover known binders from the pool of decoys. However, it is overambitious to expect any docking algorithm to recover all the known binders and none of the decoys. To quantify the {\color{red}{accuracy}} of the docking algorithms and the parameters used, an enrichment factor (EF)\cite{bib4} is used, which is computed using Eq \textbf{\ref{eq:EF-recover}}.

\begin{equation} 
  \label{eq:EF-recover}
   \ EF = \frac{\frac{ligands_{(recovered)}}{N_{(recovered)}}}{\frac{ligands_{(total)}}{N_{(total)}}} 
\end{equation}

Here, $ligands_{(total)}$ and $ligands_{(recovered)}$ denote the total number of known binders and binders recovered, respectively, and $N_{(total)}$ and $N_{(recovered)}$ denote the total number of decoys and recovered decoys, respectively. From the formula, we can infer that a higher value suggests an effective docking algorithm and parameters thereof. A comprehensive discussion on search algorithms and scoring functions is beyond the scope of this article, and the readers are advised to read the listed papers in the literature \cite{bib5, bib6}. However, we enlist the scoring functions with their succinct descriptions. 

\subsubsection*{Force field-based scoring function:}

Force fields\cite{bib28} are mathematical frameworks that describe the potential energy of a molecular system as a function of the positions of its atoms. In other words, they describe chemical information in the numerical information that is more suitable for computations. To do so, various components, each representing a specific type of interaction between atoms. Each component is characterized by its own parameters, such as the strength of the interaction and the range of distances ($r$) over which it is effective.
A force field typically consists of:
\begin{enumerate}
    \item Bonded Interactions
        \begin{enumerate}
	       \item Bond stretching
	       \item Angle bending
	       \item Torsional interactions
	       \item Improper torsions
        \end{enumerate}
    \item Non-Bonded Interactions
        \begin{enumerate}
	       \item Van der Waals forces
	       \item Electrostatics
        \end{enumerate}
\end{enumerate}

The bond stretching and angle bending are generally considered hard terms because the energy required to bring about a small change is generally very large. Moreover, to improve the computational speed, these are treated as “fixed” or “constant”. The torsional and improper terms are changed to bring about a conformational change by the conformational search algorithm. Owing to such changes, the non-bonded interactions vary as different conformations lead to different atomic interactions, which is show in Eq \textbf{\ref{eq:ffscore}} as a docking score.
    \begin{equation} 
    \label{eq:ffscore}
    \    E = \sum_{i} \sum_{j} \left(\frac{A_{ij}}{r_{ij}^{12}} - \frac{B_{ij}}{r_{ij}^{6}} + \frac{q_{i}q_{j}}{\epsilon(r_{ij}) r_{ij}}\right) \ \end{equation}
Where $r_{ij}$ is the atomic distance between protein atom i and ligand atom j, $A_{ij}$ and $B_{ij}$ are vdW repulsive and attractive parameters, respectively. $q_{i}$ and $q_{j}$ are atomic partial charges on atoms i and j. The first two terms estimate the van der Waals interactions using Lennard-Jones 12–6 potentials, which are commonly used in most force fields. The last term estimates the electrostatic component of the interaction, with solvent effects implicitly defined using the term $\epsilon(r_{ij})$, which is the distance-dependent dielectric constant. It is known that incorporating solvation effects is far from perfect, and without it, the electrostatics will be overestimated for the charged ligands owing to the absence or incorrect dielectric effect seen in the solvated systems. There are several approximations introduced such as Generalised-Born (GB)\cite{bib29} and Poisson-Boltzmann (PB) models \cite{bib30, bib31}, which are popular and computationally efficient methods that estimate solvent effects.

\subsubsection*{Empirical scoring function:}
\begin{equation} 
    \label{eq:empscore}
    \   {\Delta}G = {\sum_{i}} W_{i} \cdot {\Delta}G_{vdW} + {\sum_{j}} W_{j} \cdot {\Delta}G_{ele}  +  \ 
            {\sum_{k}} W_{solv} \cdot {\Delta}G_{solv} + . . . \
\end{equation}

Where ${\Delta}G_{vdW}$, ${\Delta}G_{ele}$, and ${\Delta}G_{solv}$, depict the free energy contribution of van der Waals', electrostatic, solvation/desolvation, etc. ${\sum_{i}} W_{i}$, ${\sum_{j}} W_{j}$, and ${\sum_{k}} W_{solv}$ are the corresponding coefficients that are adjusted to fit the experimental binding affinity data. The final score is obtained as the sum of all the energy components. Different scoring functions in this class include or exclude certain components, while others use a different fitting method that best fits the binding affinity data.
 
\subsubsection*{Knowledge-based scoring function:}
\begin{equation} 
    \label{eq:knowscore}
    \  w(r) = -k_{B}T[\frac{\rho(r)}{\rho^*(r)}] \  
\end{equation}
 
Where $w(r)$ is the final score which is dependent on the atomic distances $r$. $k_{B}$ is the Boltzmann constant, $T$ is the absolute temperature in $K$. $\rho(r)$ is the frequency of protein-ligand atom pairs at distance $r$, and $\rho^*(r)$ is the frequency of protein-ligand atom pairs in the reference state at distance $r$. The reference state is generally the frequency that a particular atom pair is observed in the experimental structures. The type of interaction (intermolecular atom pairs) that appears more frequently in the reference state will be scored higher. Therefore, the ratio $\frac{\rho(r)}{\rho^*(r)}$ can be considered as a pair distribution function.
 
\subsubsection*{Consensus scoring function:}
 \begin{equation} 
    \label{eq:scscore}
    \   {\Delta}G =  W_{1} \cdot {scoring function}_1  +  W_{2} \cdot {scoring function}_2  +  
                     W_{3} \cdot {scoring function}_3  + . . . \
\end{equation}

Where $W_{1,2,3,..}$ are weights assigned to each of the ${scoring function}$ 1, 2, 3, etc. These weights may not be directly linked to the atomic physicochemical properties. As we have seen so far, each type of scoring function comes with its own set of advantages and disadvantages. There have been many attempts to use more than one type of scoring function simultaneously to improve their overall scoring accuracy. An ideal consensus scoring function will use one or more scoring functions that complement the pros and cons. This is expected to minimise the cons and amplify the pros of each scoring function. However, this is far from reality, and often times, new limitations are observed.

Few questions that frequently arise, especially among the beginners, is \textit{what are generally the units of the docking score?} To address this question pertaining to the units, the units are generally units of energy, expresses as "kcal/mol" or "kJ/mol". However, it is not unusual to observe scoring functions that yield unitless docking scores. Irrespective, of the units of the docking score, a generally accepted convention is such that a more negative score indicates a better binding score. We would also point out that not always a more negative score indicates “correct” binding conformation (also referred to as “pose”).  
\textit{Are docking scores obtained from different scoring functions directly comparable?} This is rather difficult to answer holistically because scoring functions could be comparable as they share common philosophy of development or datasets. In our experience, we have always avoided any head-to-head comparison of scores, nevertheless, we compare the rank-ordering given by different scoring functions\cite{bib27}.  


